\section{Fully Homomorphic Encryption}


\subsection{Homomorphic Property}

The defining property of FHE is that it supports both addition and
multiplication on ciphertexts. Formally, an FHE scheme
$(\mathsf{KeyGen}, \mathsf{Enc}, \mathsf{Dec}, \mathsf{Eval})$
satisfies~\cite{gentry2009}:

\begin{theorem}[Homomorphic Correctness]
For any circuit $C$ of depth $d$, keys $(pk, sk) \leftarrow
\mathsf{KeyGen}(1^\lambda)$, and plaintexts $x, y$:
\begin{align}
\mathsf{Dec}(sk, \mathsf{Eval}(pk, +, \mathsf{Enc}(pk,x),
  \mathsf{Enc}(pk,y))) &= x + y \\
\mathsf{Dec}(sk, \mathsf{Eval}(pk, \times, \mathsf{Enc}(pk,x),
  \mathsf{Enc}(pk,y))) &= x \cdot y
\end{align}
provided the accumulated noise does not exceed the modulus bound.
\end{theorem}

The noise growth per operation is the fundamental obstacle. For
LWE-based schemes~\cite{gsw2013}, multiplication increases noise
quadratically. Gentry's bootstrapping technique~\cite{gentry2009}
resets noise by homomorphically evaluating the decryption circuit,
enabling unbounded computation depth at the cost of one bootstrapping
operation per gate.

TFHE~\cite{tfhe2016} achieves gate-by-gate bootstrapping: each
boolean gate evaluation simultaneously computes the gate and refreshes
the ciphertext noise in a single operation, avoiding noise
accumulation entirely.

\subsection{TFHE Gate-by-Gate Bootstrapping}

Each boolean gate $G \in \{\text{AND}, \text{OR}, \text{XOR}, \ldots\}$
is evaluated via a single bootstrapping operation that simultaneously
computes the gate and refreshes the noise:

\[
\text{Bootstrap}(c, G) = \text{BlindRotate}(c) \cdot G_{\text{LUT}}
\]

where $G_{\text{LUT}}$ encodes the gate's truth table as a test
polynomial in $R_q$.

\subsection{Encrypted CRDT Merge}

CRDT merge under FHE enables conflict-free replication without
decryption. The LWW-Register merge circuit compares encrypted
timestamps MSB-to-LSB:

\begin{algorithmic}[1]
\State $\mathit{bGtA} \gets \bot$; $\mathit{eqSoFar} \gets \top$
\For{$i = \text{MSB}$ \textbf{downto} $0$}
  \State $\mathit{bitGt} \gets \text{ANDNY}(a_i, b_i)$
  \State $\mathit{bitEq} \gets \text{XNOR}(a_i, b_i)$
  \State $\mathit{bGtA} \gets \text{OR}(\mathit{bGtA},
         \text{AND}(\mathit{eqSoFar}, \mathit{bitGt}))$
  \State $\mathit{eqSoFar} \gets \text{AND}(\mathit{eqSoFar},
         \mathit{bitEq})$
\EndFor
\For{each bit $i$}
  \State $\mathit{result}_i \gets \text{MUX}(\mathit{bGtA}, b_i, a_i)$
\EndFor
\end{algorithmic}

Gate count: $O(\text{bitsTS})$ comparisons + $O(\text{bitsVal} +
\text{bitsTS})$ MUX selections. At PN10QP27 with 8-bit timestamps
and 8-bit values: $\sim$40 bootstraps $\approx$ 4 seconds.

\subsection{Performance Budget}

\begin{center}
\begin{tabular}{lrrr}
\toprule
\textbf{Operation} & \textbf{Gates} & \textbf{Time (PN10QP27)} & \textbf{Ciphertext} \\
\midrule
LWW Merge (8+8 bit) & 40 & 4\,s & 136\,KB \\
GCounter Max (32 bit/node) & 128 & 13\,s & 544\,KB \\
OR-Set Merge & 0 (tag union) & $<$1\,ms & tags only \\
\bottomrule
\end{tabular}
\end{center}

%% ============================================================
